Fig.~\ref{fig:fp-protocol} summarises the protocol. Each method fingerprints a set of registered assets and a larger set of distractors, which together form the registry. Registered assets are then transformed or partially edited and used as queries; media that were never registered go through the same transformations and serve as negative queries. Every query is searched exhaustively against the registry, and a threshold fixed on one half of the negatives decides whether the best match is reported. All randomness is seeded, so every method sees byte-identical queries.

\begin{figure*}[!t]
\centering
% Benchmark protocol schematic (Fig. 1). Drawn in TikZ so text matches the document.
\begin{tikzpicture}[
  font=\sffamily\footnotesize,
  box/.style={draw=black!70, line width=0.5pt, rounded corners=1.5pt, align=center, minimum height=9mm,
              text width=30mm, inner xsep=2pt, fill=white},
  data/.style={box, fill=black!5},
  metric/.style={box, draw=none, fill=blue!8, minimum height=10mm, text width=33mm},
  arr/.style={-{Stealth[length=1.8mm]}, line width=0.5pt, draw=black!70},
  lab/.style={font=\sffamily\scriptsize, text=black!60},
]
% lane 1: registration
\node[data] (reg) at (0,0) {Registered assets\\\scriptsize images, audio, video};
\node[box] (wm) at (3.85,0) {Optional watermark\\\scriptsize TrustMark, AudioSeal,\\[-2pt]\scriptsize Video Seal};
\node[box] (fp) at (7.7,0) {Fingerprint\\\scriptsize \val{fp:count/methods} methods};
\node[data] (idx) at (11.55,0) {Registry index\\\scriptsize assets + distractors};
\draw[arr] (reg) -- (wm);
\draw[arr] (wm) -- (fp);
\draw[arr] (fp) -- (idx);
\node[lab, anchor=south] at (5.775,0.72) {registration};

% lane 2: queries
\node[data] (neg) at (0,-1.9) {Never-registered\\\scriptsize media (negatives)};
\node[box] (att) at (3.85,-1.9) {Transform or edit\\\scriptsize \val{fp:count/attacks} attacks, 6 edits};
\node[box] (qfp) at (7.7,-1.9) {Fingerprint\\\scriptsize query};
\node[box] (srch) at (11.55,-1.9) {Exact search\\\scriptsize top-50, threshold};
\draw[arr] (reg.south east) -- node[lab, right, pos=0.45, inner sep=3pt] {positives} (att.north west);
\draw[arr] (neg) -- node[lab, below, inner sep=2pt] {negatives} (att);
\draw[arr] (att) -- (qfp);
\draw[arr] (qfp) -- (srch);
\draw[arr] (idx) -- (srch);
\node[lab, anchor=north] at (5.775,-2.55) {queries};

% lane 3: what is measured
\node[metric] (m1) at (0,-3.75) {\textbf{Robustness}\\\scriptsize R@1, $\mu$AP, TPR at calibrated FPR};
\node[metric] (m2) at (3.85,-3.75) {\textbf{Localization}\\\scriptsize edited region; time offset};
\node[metric] (m3) at (7.7,-3.75) {\textbf{Security}\\\scriptsize evasion, forgery, inversion};
\node[metric] (m4) at (11.55,-3.75) {\textbf{Complementarity}\\\scriptsize watermark key vs.\ fingerprint};
\draw[densely dashed, black!50, line width=0.5pt] ($(m1.north west)+(-0.1,0.18)$) -- ($(m4.north east)+(0.1,0.18)$);
\end{tikzpicture}

\caption{Benchmark protocol. \val{fp:count/methods} fingerprinting methods are applied to registered assets and distractors (registry) and to queries derived from registered assets (positives) or from never-registered media (negatives) under \val{fp:count/attacks} transformations and six kinds of partial edit. Thresholds are fixed on half of the negatives and evaluated on the other half. The same queries also drive the localization, security and watermark-complementarity measurements.}
\label{fig:fp-protocol}
\end{figure*}

\subsubsection{Methods Under Test}
We included every fingerprinting method we could find that is distributed with an implementation and, where learned, with weights, and that targets natural images, music and speech, or natural video (Table~\ref{tab:fp-methods}). We recorded the licence of the code and, separately, of the weights and of the data they were trained on, because the three differ in several projects. A method is \emph{product-eligible} (tier~P) when code and weights carry a permissive licence (MIT, BSD, Apache-2.0 or CC0) and the weights were not trained on data restricted to non-commercial use. The strongest learned copy detectors are not: SSCD and the ISC21 winner are MIT-licensed code, but their weights carry no licence of their own and were trained on DISC21, whose repository is CC BY-NC 4.0; DINOHash has no licence file; and the neural music fingerprint NMFP is GPL-3.0. We measured them as reference points (tier~R, marked \reftier), as the watermarking track does for its unlicensed checkpoints (Section~\ref{sec:wm-method}). We excluded Apple NeuralHash (weights extracted from a device, no licence), PhotoDNA (licensed only for child-safety use), Adobe's image comparator network (no public code), SuperPoint and TruFor (non-commercial weights), and Olaf and Panako, which we did not run: they are AGPL-3.0, whose network-use clause rules them out for a hosted provenance service, and Chromaprint and audfprint already represent their landmark and sub-band designs.

\paragraph{Images.} Twenty-one descriptors: PDQ \cite{meta2019pdq}, with and without matching the query's eight dihedral variants; aHash, dHash, pHash (64 and 256 bits) and wHash from imagehash \cite{buchner_imagehash}; BlockMean, Marr--Hildreth and ColorMoment from OpenCV \cite{opencv_imghash}; Blockhash \cite{commonsmachinery_blockhash}; the ISCC Image-Code at 64 bits (the ISO default) and 256 bits \cite{iso24138}; DINOv2-S and DINOv2-B class tokens \cite{oquab2024dinov2}; a keyed 256-bit sign projection of DINOv2-S (below); OpenCLIP ViT-B/32 \cite{cherti2023openclip}; SSCD with ResNet-50 and ResNeXt-101 backbones \cite{pizzi2022sscd}; the ISC21 descriptor-track winner \cite{yokoo2021isc}; and DINOHash \cite{singhi2025dinohash}. ISCC's experimental semantic image code \cite{iscc_sci} ran in our pilot but not at full scale: it binarises an ISC21-derived descriptor that we evaluate directly, and its reference model runs one image at a time at about 3\,s per image on a CPU core. Learned descriptors see the whole image resized to their square input without cropping, so that a crop attack cannot remove content the descriptor was never shown. Blockhash's reference implementation loops over pixels in Python; we vectorised it and verified bit-identical output on five image sizes including odd ones.

\paragraph{Keyed projection.} DINOv2-S LSH-256 multiplies the L2-normalised DINOv2-S descriptor by a $384 \times 256$ Gaussian matrix drawn from a secret seed and keeps the signs. It is a standard random-hyperplane hash, included because it is compact, product-eligible and, unlike a raw descriptor, meaningless to anyone without the key; the security experiments test that claim.

\paragraph{Audio.} Chromaprint \cite{lalinsky_chromaprint}, with candidate generation on the 20 most significant bits of each sub-fingerprint and verification by bit error rate at the voted offset (the AcoustID index design); the ISCC Audio-Code at 64 and 256 bits, computed from the whole-file Chromaprint vector; audfprint's spectral-peak landmarks \cite{ellis_audfprint}, scored by the number of hashes that agree on one time offset; LAION-CLAP \cite{wu2023clap} as a semantic embedding, mean-pooled over 10\,s windows; and NMFP \cite{araz2025nmfp}, a neural music fingerprint with 1\,s segment embeddings, run in its authors' TensorFlow environment and matched by segment voting.

\paragraph{Video.} vPDQ \cite{meta_vpdq} at one hash per second, scoring a candidate by the share of query frames with a reference frame within Hamming distance 31; TMK+PDQF \cite{meta2019pdq, poullot2015tmk}, scored by its level-2 time-kernel similarity using Meta's command-line tools; the ISCC Video-Code at 64 and 256 bits from MPEG-7 frame signatures \cite{paschalakis2012mpeg7}; and DINOv2-S and SSCD on frames sampled at 1\,fps, either mean-pooled into one clip descriptor or matched frame by frame with temporal voting, the baseline recipe of the 2023 Video Similarity Challenge \cite{pizzi2024vsc}.

\begin{table*}[!t]
\centering
\scriptsize
\caption{Methods, licences, stored size and extraction time. Bits: stored bits per registered asset (float descriptors count 32 bits per dimension; time-based systems report the mean over the registered items). Time: per image (1024 px), per second of audio or per 5\,s clip.}
\label{tab:fp-methods}
\setlength{\tabcolsep}{3pt}
\begin{tabular}{lllp{3.3cm}p{3.6cm}rr}
\toprule
Modality & Method & Class & Code licence & Weights / data & Bits & Time (ms) \\
\midrule
Image & PDQ & Perceptual hash & BSD (ThreatExchange); pdqhash binding  & n/a & 256 & 31 CPU \\
 & PDQ (dihedral) & Perceptual hash & BSD (ThreatExchange); pdqhash binding  & n/a & 256 & 31 CPU \\
 & aHash & Perceptual hash & BSD-2-Clause (imagehash) & n/a & 64.0 & 4.3 CPU \\
 & dHash & Perceptual hash & BSD-2-Clause (imagehash) & n/a & 64.0 & 4.4 CPU \\
 & pHash & Perceptual hash & BSD-2-Clause (imagehash) & n/a & 64.0 & 5.2 CPU \\
 & pHash-256 & Perceptual hash & BSD-2-Clause (imagehash) & n/a & 256 & 5.6 CPU \\
 & wHash & Perceptual hash & BSD-2-Clause (imagehash) & n/a & 64.0 & 63 CPU \\
 & BlockMean & Perceptual hash & Apache-2.0 (opencv\_contrib) & n/a & 256 & 1.0 CPU \\
 & Marr–Hildreth & Perceptual hash & Apache-2.0 (opencv\_contrib) & n/a & 576 & 5.5 CPU \\
 & ColorMoment & Perceptual hash & Apache-2.0 (opencv\_contrib) & n/a & 1344 & 3.6 CPU \\
 & Blockhash & Perceptual hash & MIT & n/a & 256 & 19 CPU \\
 & ISCC Image-64 & ISCC (ISO 24138) & Apache-2.0 & n/a & 64.0 & 11 CPU \\
 & ISCC Image-256 & ISCC (ISO 24138) & Apache-2.0 & n/a & 256 & 11 CPU \\
 & DINOv2-S & General embedding & Apache-2.0 & Apache-2.0 (HF model card) & 12288 & 2.0 GPU \\
 & DINOv2-B & General embedding & Apache-2.0 & Apache-2.0 (HF model card) & 24576 & 3.9 GPU \\
 & DINOv2-S LSH-256 & General embedding & Apache-2.0 & Apache-2.0 (HF model card) & 256 & 1.7 GPU \\
 & OpenCLIP B/32 & General embedding & MIT (open\_clip) & MIT (laion/CLIP-ViT-B-32-laion2B-s34B-b79K) & 16384 & 1.1 GPU \\
 & SSCD R50\reftier & Copy-detection trained & MIT & none stated; trained on DISC21 (CC BY-NC 4.0 & 16384 & 2.7 GPU \\
 & SSCD RX101\reftier & Copy-detection trained & MIT & none stated; trained on DISC21 (CC BY-NC 4.0 & 32768 & 4.6 GPU \\
 & ISC21 1st\reftier & Copy-detection trained & MIT & MIT repo release; trained on DISC21 (CC BY-N & 8192 & 8.3 GPU \\
 & DINOHash-96\reftier & Copy-detection trained & none (no LICENSE file) & none (no LICENSE file) & 96.0 & 3.6 GPU \\
\addlinespace
Audio & Chromaprint & Perceptual hash & MIT (chromaprint core; fpcalc binary l & n/a & 7072 & 2.0 CPU \\
 & audfprint & Perceptual hash & MIT & n/a & 40397 & 1.4 CPU \\
 & ISCC Audio-64 & ISCC (ISO 24138) & Apache-2.0 & n/a & 64.0 & 2.0 CPU \\
 & ISCC Audio-256 & ISCC (ISO 24138) & Apache-2.0 & n/a & 256 & 2.0 CPU \\
 & CLAP & General embedding & CC0-1.0 & CC0-1.0 (LAION-AI/CLAP release) & 16384 & 1.1 GPU \\
 & NMFP\reftier & Copy-detection trained & GPL-3.0 & Zenodo 15719945 (licence as repo, GPL-3.0);  & 239108 & -- \\
\addlinespace
Video & vPDQ & Perceptual hash & BSD (ThreatExchange vpdq) & n/a & 1258 & 132 CPU \\
 & TMK+PDQF & Perceptual hash & BSD (ThreatExchange tmk) & n/a & 2106752 & 51 CPU \\
 & ISCC Video-64 & ISCC (ISO 24138) & Apache-2.0 & n/a & 64.0 & 54 CPU \\
 & ISCC Video-256 & ISCC (ISO 24138) & Apache-2.0 & n/a & 256 & 54 CPU \\
 & DINOv2-S frames & General embedding & Apache-2.0 & Apache-2.0 (HF model card) & 61440 & 136 GPU \\
 & DINOv2-S mean & General embedding & Apache-2.0 & Apache-2.0 (HF model card) & 12288 & 136 GPU \\
 & SSCD frames\reftier & Copy-detection trained & MIT & none stated; trained on DISC21 (CC BY-NC 4.0 & 81920 & 137 GPU \\
 & SSCD mean\reftier & Copy-detection trained & MIT & none stated; trained on DISC21 (CC BY-NC 4.0 & 16384 & 137 GPU \\
\bottomrule
\end{tabular}
\par\vspace{2pt}\parbox{\linewidth}{\raggedright\scriptsize \reftier~Reference tier: weights without a licence, trained on non-commercial data, or copyleft code; measured for comparison, not product-eligible.}
\end{table*}


\subsubsection{Corpora}
\emph{Images.} Amazon Berkeley Objects (ABO, CC BY 4.0) \cite{collins2022abo} supplies product photographs from real marketplace listings, the setting in which copied and relabelled product images circulate. We sampled \val{fp:corpus/abo/reg} registered originals (one main image per product, shorter side at least 800\,px), stored at a longer side of 1024\,px as the asset a provenance service would register, and \val{fp:corpus/abo/dist} distractors (one image per other product). Queries come from \val{fp:corpus/abo/pos} registered originals; negative queries from \val{fp:corpus/abo/neg} images of products that are neither registered nor distractors; and hard negatives from \val{fp:corpus/abo/hard} \emph{other} photographs of registered products. ABO reuses photographs across listings and colour variants, so before scoring we flagged pairs that are copies of each other under a rule fixed in advance on unattacked originals (PDQ distance at most 16 bits, SSCD cosine at least 0.80, or DINOv2-B cosine at least 0.95). The rule removed \val{fp:dedup/n_dist_flagged} distractors and \val{fp:dedup/n_neg_flagged} negative sources, and marked \val{fp:dedup/n_hard_near_duplicate} of \val{fp:dedup/n_hard} hard negatives as recoloured versions of the registered photograph, which we report separately. As a second, literature-comparable set we fetched a subset of the DISC21 development split \cite{douze2021isc}: \val{fp:corpus/disc/gt} ground-truth query--reference pairs, \val{fp:corpus/disc/negq} queries without a reference, and \val{fp:corpus/disc/refs} references in total. DISC21 is licensed for non-commercial research; we use it for measurement only.

\emph{Audio.} \val{fp:corpus/audio/music} music tracks from the FMA large set \cite{defferrard2017fma}, restricted to tracks released under CC BY, CC BY-SA or CC0 and excluding the FMA-medium subset, and \val{fp:corpus/audio/speech} LibriSpeech utterances of at least 8\,s \cite{panayotov2015librispeech}, all at 22.05\,kHz mono. Music clips are 30\,s. Of these, \val{fp:corpus/audio/reg} are registered, \val{fp:corpus/audio/dist} are distractors and \val{fp:corpus/audio/neg} are negative sources.

\emph{Video.} Five-second clips cut every 5\,s from four Blender open films \cite{blender2010sintel, blender2008bbb, blender2012tos, blender2006ed}, skipping titles and end credits, and the first 5\,s of \val{fp:corpus/video/xiph} Xiph test sequences \cite{xiphderf}, scaled to 1280\,px wide at 25\,fps and stored as H.264 at CRF~12: \val{fp:corpus/video/reg} registered, \val{fp:corpus/video/dist} distractor and \val{fp:corpus/video/neg} negative clips. Clips from the same film share style and characters, which makes the distractors deliberately hard.

\subsubsection{Transformations}
\emph{Images} (\val{fp:count/image/attacks}, grouped into families in Fig.~\ref{fig:fp-image-heatmap}): JPEG at quality 75, 50 and 30 and WebP at 50; downscaling to 0.5 and 0.25 and a 0.75 aspect change; centred crops keeping 90, 70, 50 and 30\% of each side and an off-centre 50\% crop; rotation by 5, 15 and 90 degrees, horizontal flip, perspective, skew and 25\% padding; brightness, contrast, saturation, gamma, greyscale, blur, Gaussian noise and pixelation; text, emoji (30\% of the image), meme caption, screenshot and overlay onto another photograph at half size, following AugLy \cite{papakipos2022augly}; and two chains, a social-media chain (downscale, JPEG~70, 90\% crop) and a re-capture simulation (perspective, defocus, sensor noise, tone curve, JPEG~80). \emph{Audio} (\val{fp:count/audio/attacks}): MP3 and AAC at 64\,kb/s, Opus at 24\,kb/s, resampling through 8\,kHz; white noise at 20 and 10\,dB SNR and babble from other music at 10 and 0\,dB; telephone band, 2\,kHz low-pass, dynamic compression and reverberation; pitch $+2$ semitones, tempo 0.9 and 1.1 and speed 1.05; excerpts of 10 and 5\,s at a random offset, an excerpt with MP3 and noise, and a re-recording simulation. Unlike the watermarking track, codec outputs are not re-aligned: a fingerprint must cope with codec delay itself. \emph{Video} (\val{fp:count/video/attacks}): H.264 at CRF~28 and 35, H.265 at CRF~32; downscaling to 0.5 and 0.25; 75\% and 50\% crops; flip, 90 and 5 degree rotation, letterboxing; greyscale, brightness, contrast, blur and noise; text caption, logo, and picture-in-picture at half size on another clip; 15\,fps, 1.25$\times$ speed, a 2.5\,s excerpt, the same 2.5\,s inserted into 5\,s of another clip, and a screen-recording simulation.

\subsubsection{Retrieval Metrics and Operating Points}
Search is exact (flat index, top~50): Hamming similarity for binary codes, cosine for L2-normalised descriptors, and each audio or video system's own score. We report Recall@1 and the micro-average precision ($\mu$AP) of the ISC21 protocol \cite{douze2021isc}, both threshold-free. A registry, however, must decide, and a false decision binds content to someone else's manifest. We therefore fix thresholds on the \emph{calibration} half of the negative sources (all their transformed copies) and report on the other half. Two kinds of target are used. The query-level target fixes the rate at which a never-registered query's best match passes at the benchmark's registry size (1\%). The pair-level target fixes the rate at which a (negative query, reference) pair passes. For a registry of $N$ assets and an average pair-level false-match probability $p$, the expected number of false matches per negative query is $Np$; the probability of at least one false match is at most $\min(1,Np)$, without requiring independent reference comparisons. Our headline target is a pair-level rate of $10^{-7}$: at the benchmark's \val{fp:count/image/nref} references, this corresponds to about one expected false match per hundred negative queries in aggregate. This calculation is not a measured query-level false-binding rate. Table~\ref{tab:fp-ci} reports held-out query-level rates at the separately calibrated query-level threshold. We also report $10^{-8}$, the smallest rate our \val{fp:count/image/neg_cal} calibration queries resolve. A positive query counts as detected only if its best match is the true asset and passes the threshold. A target is reported only if the calibration set resolves it, that is, if it implies at least five false pairs (or queries) among the calibration negatives; otherwise the threshold would merely sit above every negative observed. This is a reporting-resolution rule, not evidence that the comparisons are independent or that the tail probability is estimated precisely. The queries are not independent: every registered or negative source contributes one query per transformation, and the transformed copies of one source tend to succeed or fail together. Intervals for these retrieval metrics therefore come from a bootstrap over sources (2{,}000 replicates): registered sources are drawn with replacement together with all their queries, calibration negative sources likewise, with the threshold recomputed from the drawn calibration set in every replicate, and held-out negative sources for the out-of-sample false-match rate. Calibration and held-out negatives come from disjoint sources. Planned paired comparisons use a sign-flip test on each source's difference in detections between the two methods, with Holm adjustment \cite{holm1979}, and report a bootstrap interval of the difference in detection rate; both methods are resampled with the same sources. The sign-flip test and the verification experiment (Table~\ref{tab:fp-verify}) hold the thresholds at their full-sample values, so their uncertainty excludes threshold estimation; the bootstrap interval of a paired difference recomputes the thresholds in every replicate.

\subsubsection{Partial Edits and Localization}
Partial edits are generated from \val{fp:edit/n_sources} registered originals, with the edited region known exactly: a smooth blob of 1, 3, 10 or 25\% of the image, placed where the image has texture. The edit is a splice from another photograph, a copy-move within the image, classical (Telea) or learned (LaMa \cite{suvorov2022lama}) inpainting of the region, a hue rotation, or replacement by a flat label with new text, followed by no post-processing, JPEG~75, downscaling with JPEG~80, or a 90\% crop. A balanced design gives \val{fp:edit/n_edits} edited images. We first ask whether each fingerprint still retrieves the original at its calibrated threshold. We then localize the edit against the registered original: the query is registered onto the original with SIFT and RANSAC \cite{lowe2004sift, fischler1981ransac} or with DISK and LightGlue \cite{tyszkiewicz2020disk, lindenberger2023lightglue}, and compared through six difference maps: blurred pixel difference, SSIM \cite{wang2004ssim}, LPIPS \cite{zhang2018lpips}, DINOv2-S patch distance, PDQ over an $8 \times 8$ grid of tiles, and a fusion of the five. Map thresholds are fixed at a pixel false-positive rate of 1\% on unedited originals put through the same post-processing, on half of the sources; the other half is evaluated. For audio and video we measure temporal localization: how accurately each system recovers the offset of an excerpt within its source, and which interval of an inserted clip matches.

\subsubsection{Security}
We consider an adversary who knows the fingerprinting method and wants either to \emph{evade} it (republish a registered asset so that it no longer matches) or to \emph{forge} a match (make an unrelated image bind to a registered asset, for example to borrow its credentials). Evasion uses projected gradient descent \cite{madry2018pgd} under an $L_\infty$ budget of 2, 4, 8 or 16 grey levels: white-box on the learned descriptors, through differentiable re-implementations of PDQ, pHash, aHash, dHash and the ISCC Image-Code for those hashes, and by transfer from an ensemble of DINOv2-S, SSCD and OpenCLIP for all methods. Success is always judged by the real extractor on the 8-bit result. Image and audio attacks were judged at the method's pair-level $10^{-6}$ threshold, and video attacks at its query-level 1\% threshold, because the video calibration set does not resolve $10^{-6}$ (Section~\ref{sec:fp-res-video}). For images, $10^{-6}$ is a looser operating point than the headline: an evasion must push the score below it, which makes the reported evasion rates conservative, and we re-score white-box evasions at the $10^{-7}$ threshold from their final scores. Forgery success at $10^{-6}$ is correspondingly an upper bound on success at $10^{-7}$. Forgery pushes a negative image towards a registered target and succeeds only if the target becomes the top match in the full registry above threshold. Because a soft-binding value is written into the manifest, we also test \emph{inversion}: a decoder trained on distractor descriptors reconstructs a 64$\times$64 thumbnail from a registered asset's descriptor, and we measure whether the reconstruction identifies the asset among the others. For audio and video we apply the same attacks to the learned systems and measure, for every method, the smallest pitch, tempo, speed or crop change that breaks it.

\subsubsection{Watermark Complementarity}
The registered assets of the three modalities were also watermarked with MIT-licensed models from the watermarking track: TrustMark~Q at its default strength for images \cite{bui2025trustmark}, AudioSeal \cite{sanroman2024audioseal} and Video Seal \cite{fernandez2024videoseal}. The marked copies went through the same transformations. For every query we record whether the watermark satisfies its method-specific success criterion, and whether each fingerprint identifies the asset. The three watermarks decide this differently, and the decisions relate differently to recovering a registry identifier. TrustMark transmits a 100-bit codeword; in the BCH\_5 mode we use it carries a 61-bit data payload, 35 parity bits and 4 version bits. Its BCH decoder either corrects the codeword or reports failure; the key counts as recovered when the decoded 61-bit payload equals the asset's key exactly, which is identification: the decoded payload itself names the registry entry. AudioSeal (the base model) carries only 16 bits and has no error correction. Its detection score is the fraction of audio frames classified as watermarked; we require at least half of the frames, a more permissive rule than the 0.8 of the watermarking track's blind detector (Section~\ref{sec:wm-fp-method}), because here detection counts only together with an exact match of all 16 bits to the asset's key. No AudioSeal decode of the watermarking track's unmarked audio exceeded \val{wm:audioseal/unmarked_max_score}, so its blind false-positive count is the same under either rule. A recovered key identifies an asset only if no other asset holds it. The space contains only $2^{16}$ keys, so a registry larger than that cannot give every asset a distinct AudioSeal key; our keys were drawn pseudo-randomly per asset without a uniqueness check, and the \val{fp:wm/audio/n_assets} registered assets received \val{fp:wm/audio/nkeys} distinct keys, so \val{fp:wm/audio/n_colliding} assets share one key. Video Seal carries 256 bits without error correction, and exact recovery of all 256 bits is rare after compression; we count the query as passing expected-key verification when at least $k$ of the 256 decoded bits agree with the asset's key, with $k$ the smallest value for which a random 256-bit word agrees with a given key with probability at most $10^{-6}$. For positives this is verification against the expected key, not a test of unique identification among all registered keys. A registry reader without that expectation would compare the decoded bits with every registered key; on the never-registered negatives we did exactly that, so a false Video Seal key there means agreement with \emph{any} of the registered keys. Under the independent-fair-bit null used to set $k$, the union bound gives a per-query false-match probability of at most $\min(1,N \times 10^{-6})$ for $N$ keys. This analytic bound is conditional on that null and does not replace the measured false-match rate. Testing all keys on negative queries does not establish unique identification on positive queries. For AudioSeal negatives we report both the detector's false detections and detections whose message equals a registered key. We also measure how far the watermark itself moves each fingerprint, and whether the registry should hold the fingerprint of the original or of the marked asset.

\subsubsection{Second-Stage Verification}
To test whether checking a retrieved candidate against the stored original removes the false matches that a global descriptor lets through, we ran SIFT with RANSAC between every image query (positive and negative) and the registered original of its top-1 candidate under DINOv2-S and PDQ, at the localization track's working size, and recorded the number of geometric inliers. A match is accepted when the retrieval score passes a threshold and the inlier count reaches $k$; both are fixed on the calibration negatives and evaluated on the held-out ones.

\subsubsection{Hosts and Reproducibility}
Learned descriptors were extracted on one NVIDIA A10G (AWS g5.2xlarge, AMD EPYC 7R32); CPU descriptors single-threaded on the same processor model (an AWS c5a.16xlarge was added to spread the CPU work). Every source tree, checkpoint and corpus file is pinned by commit or SHA-256. Because the watermarking track found host-dependent decoder outputs (Section~\ref{sec:wm-res-stability}), we fingerprinted the same 200 PNG files on the GPU host and on an Apple M5 Pro laptop and compared the descriptors bit by bit.

The retrieval tracks and the watermark combination were run twice, on 2026-09-26 and again on 2026-09-28 on fresh hosts of the same types, because the first run did not keep the per-query scores that source-level intervals require. The corpus was rebuilt from its sources and matched the pinned SHA-256 of every image, audio and video item. The image track reproduced the first run exactly (every pooled rate to within $10^{-6}$), and the video track to within two of 3{,}000 queries. The audio track did not reproduce exactly, for two reasons that we found and fixed: the background music of the two babble transformations had been drawn from whichever surplus FMA downloads were present, which differed between runs (the pool is now pinned by name and hash), and CLAP was extracted on the CPU host in the second run instead of the GPU. All reported numbers come from the second run.
